\begin{table}
\small
\centering
\begin{tabular}{lp{.75\textwidth}}
Step & Prompt \\ \hline
Intermediate Rep & "Given the input text, extract the document title and authors. For each section in the given input text, extract the most important sentences. Format the output using the following JSON template:\textbackslash{}n \textless{}JSON STRUCTURE\textgreater \textbackslash{}n\textbackslash{}n ~Input: \textless{}INPUT DOC\textgreater\textbackslash{}n\textbackslash{}n Output:" \\
\hdashline[0.5pt/3pt] 
Templatic View & "Summarize the following input in a \textless{}TEMPLATE TYPE\textgreater ~style.~Style parameters: \textless{}STYLE PARAMS\textgreater~Format the output document as a latex file: \textbackslash{}n\textbackslash{}nInput: \textless{}INPUT DOC\textgreater \textbackslash{}n\textbackslash{}n Output:"
\end{tabular}
    \caption[Prompts used to generate the intermediate representation and final LaTeX document.]{Prompts used to generate the intermediate representation and final LaTeX document. The JSON structure is pictured in Figure~\ref{fig:sure-template} and exact style parameters used can be found in Table~\ref{tbl:style-params}.} %\isabel{move to appendix if space is needed}}
    \label{tbl:prompts}
\end{table}
\begin{figure}[t]
    \centering
\includegraphics[width=.8\columnwidth]{research-chapters/3-chapter/images/sure_template.pdf}
    \caption[Abbreviated template of the intermediate representation.]{Abbreviated template of the intermediate representation provided to the prompts in Table~\ref{tbl:prompts}. We note that the specific structure provided to the prompt is not inherent to our method, and a different structure could be provided depending on the input document and domain.}
    \label{fig:sure-template}
\end{figure}
% \input{tables/prompts-style}
\section{Unified LLM-powered Generation of Templatic Views}\label{sec:3-chapter-method}
The most straightforward way to transform documents between templatic views using LLMs is to simply prompt the system to generate the target view given the input. However, similar to chain of thought prompting~\cite{Wei2022ChainOT}, we hypothesize that first generating a structured, intermediate representation of an input document and then reasoning over that representation will result in better generations than directly prompting the model. Our goal is to evaluate the capabilities of LLMs to generate long, structured documents, and experiment with how structured prompting can improve performance. We experiment with a simple general two-step process: first generate an intermediate representation, then generate the templatic view. These steps are described in greater detail below, and the process is visualized in Figure~\ref{fig:method}.


\paragraph{Intermediate Representation Generation.} In this chapter, we set the intermediate representation to be a JSON consisting of a structured layout of the most important parts of the input. We provide the input document to the model along with a template of the representation and prompt it to extract the most important information from the input document, and format it in the given JSON structure. The prompts can be found in Table~\ref{tbl:prompts} and the JSON structure can be found in Figure~\ref{fig:sure-template}. 
While our experiments use a JSON intermediate representation, note that other formats that provide structure to the input text could be employed (e.g. XML or Markdown). Rather than trying to optimize for the best representation format, our goal is to show that this chain of extraction approach along with structured augmentation to prompts can aid the quality of generations from LLMs. We leave exploration of different formats and other prompt optimization to future work. 

\begin{table}[]
\small
\centering
\begin{tabular}{lp{.8\textwidth}}
 Template & Style Params \\ \hline
Slides & ``Slides should include a title page. Following  slides should contain an informative slide  title and short, concise bullet points. Longer slides should be broken up into multiple slides.''
\\ \hdashline[0.5pt/3pt]
Posters & ``Posters should include a title section at the top. Each panel should include a heading and short, concise bullet points of the most important take-aways from that section.''
\\ \hdashline[0.5pt/3pt]
Blogs & ``Blogs should include paragraphs introducing the topic, a summary of the input document, and important takeaways. The blog should be more readable to a general audience than the input document.''
\end{tabular}
\caption[Style parameters.]{Style parameters provided to the prompts in Table~\ref{tbl:prompts}. Style parameters allow the user to specify template style with minimal prompt engineering.}
\label{tbl:style-params}
\end{table}

\paragraph{Templatic View Generation.} We then feed the generated representation as input back into the LLM, prompting the model to generate the final output document, represented as a LaTeX document. For each templatic view, the prompt to generate the final LaTeX document takes a short description of the desired output, which we refer to as a style parameter. 
The style parameters we used can be found in Table~\ref{tbl:style-params}. Style parameters make our method adaptable to new templatic views; the user only needs to write a short description of the template style. Both the generation of the intermediate representations and the final documents require little to no prompt engineering. 


